\section{Literature review}\label{sec:literature}

In this review, we aim to connect three literatures in sequence: textual
analysis in finance, the
geometry of probability measures through optimal transport, and
characteristic-based and spatial representations of systematic risk.
Together, these strands provide the measurement object, the distance used to
compare it, and the economic second moment that the paper seeks to restrict.

Textual analysis in finance increasingly uses language models to encode
heterogeneous firm information in a common representation space.
Finance-specific models such as FinBERT improve sentiment analysis and
information extraction from financial text \citep{huang_finbert_2023}.
Building on these representations, more recent applications use
language-model outputs and embeddings to forecast returns and study how
markets process news \citep{chen_expected_2026,lopezlira_chatgpt_2026}.
These applications largely map individual documents into classifications,
scores, or expected-return predictions.
This paper addresses that measurement limit by retaining each firm's
collection of representations and relating distributional separation to
systematic covariance rather than another document-level prediction.

Optimal transport supplies the geometry for that comparison.
Word Mover's Distance treats word embeddings as empirical distributions and
compares texts by the cost of transporting probability mass through semantic
space \citep{kusner_word_embeddings_2015}.
This construction shows how a collection of language-model representations
can become a probability measure whose geometry reflects distances in the
underlying representation space.
Word Mover's Distance does not by itself link distributional separation to a
restriction on financial second moments; this paper develops that link for
systematic covariance.

The relevant geometric foundation is Wasserstein optimal transport, which
connects to but does not require the broader machinery of information geometry.
Classical information geometry studies families of probability distributions
as points on geometric manifolds and commonly measures their separation
through divergences such as Kullback--Leibler.
Wasserstein information geometry instead builds on the cost of moving
probability mass through the underlying sample space
\citep{amari_information_wasserstein_2018,khan_optimal_transport_information_2022}.
Unlike classical information geometry, which is invariant to
reparametrization, the Wasserstein metric retains the cost structure of that
sample space: transport between nearby regions is cheap, whereas transport
between distant regions is expensive.
Through Kantorovich duality, this cost has a direct economic reading as the
maximum valuation difference between two configurations of mass over price
functions that are free of spatial arbitrage.
Recent work develops Wasserstein statistical manifolds, transport-induced
natural gradients, and Wasserstein information matrices to connect
information geometry and optimal transport
\citep{li_natural_gradient_2018,panaretos_invitation_2020,li_wasserstein_information_2023}.
Those tools primarily address statistical efficiency and estimation on
probability models, whereas the present paper uses the underlying Wasserstein
transport geometry to study an economic second moment.
Our covariance argument therefore does not invoke the broader
information-geometric machinery; it asks what the optimal-transport distance
between two probability laws can imply about their feasible covariance.

The factor-model literature supplies the financial object to which that
geometry is applied: systematic covariance generated by exposure to common shocks.
In the Arbitrage Pricing Theory of \citet{ross_arbitrage_1976}, returns load
on pervasive factors and no-arbitrage restrictions discipline expected returns.
Focusing instead on second moments, \citet{chamberlain_rothschild_1983} study
approximate factor structure through the spectrum of the return covariance operator.
Characteristic-based factor-loading models treat exposures as functions of
firm descriptors
\citep{rosenberg1974extra,connor2007semiparametric,connor2012efficient}, and
the time-varying-beta literature allows those loadings to change
\citep{ghysels_stable_1998}.
Proxy-selection results further caution that proxy factors inherit
approximate pricing relations only under maintained factor structure and
additional rank conditions
\citep{reisman1992reference,shanken1992current,middleton2000deriving}.
Taken together, this literature distinguishes factor exposures, their
covariance implications, and the additional conditions required for
expected-return restrictions.
Our target is systematic covariance rather than expected returns because the
alignment of common-shock exposures determines how systematic risk aggregates
across firms without requiring a new no-arbitrage pricing claim.
Relative to these models, the paper changes the representation of firm
characteristics from a single descriptor vector to a probability law and asks
how distance between two such laws restricts the covariance their induced
exposures can support.

A complementary finance literature organizes risk through continuous spaces,
asset-indexed fields, and economically meaningful notions of proximity.
Random-field representations have been used for the term structure
\citep{kennedy_term_1994,goldstein_term_2000,santa_clara_sornette_2001}.
Related constructions model equity returns with a stochastic string
\citep{distaso_string_2024} and liquidity with an infinite-dimensional field
\citep{lototsky_liquidity}.
Continuum-of-assets models instead place the continuum over securities
\citep{werner1985,al1995decomposition,khan2001asymptotic}.
Spatial and network factor models encode economic proximity through
geographic, supplier--customer, or news co-mention links
\citep{menzly_market_2010,fernandez_spatial_2011,scherbina_economic_2013,bera_spatial_2016,kou_asset_2018,schwenkler_network_2020,ge_news-implied_2023,gao2025high}.
These approaches show that economic proximity can organize common risk.
Yet spatial and network models typically encode proximity with a discrete
relation matrix, while characteristic-based models represent each firm by a
finite descriptor or loading vector; neither formulation represents within-firm
heterogeneity as a distribution.

We instead represent each firm's text-derived characteristics as an empirical
probability measure and compare firms using quadratic Wasserstein distance.
Maintained restrictions on a common characteristic-to-exposure map transfer
bounds on that distance to latent exposure space.
There, polarization converts the minimum expected squared distance between
exposure laws into the attainable ceiling on pairwise systematic covariance.
Article-embedding geometry serves only as a proxy for the latent characteristic
laws and does not identify the exposure map or its restrictions.
The next section formalizes this bridge and derives the resulting one-sided
envelope for the covariance ceiling.
